\documentclass[12pt]{article}
\usepackage[utf8]{inputenc}
\usepackage[T1]{fontenc}
\usepackage{hyperref}
\title{Arin Language Resources for K-12 Curriculum Development}
\author{Compiled by Claude Sonnet 5.5 (claude-sonnet-5-5) for Andr\'as Kornai}
\date{}
\begin{document}
\maketitle

\section{Introduction}
Arin (ISO 639-3: xrn; Glottocode arin1243) is an extinct Yeniseian language of Siberia. It was spoken along the Yenisei River between Yeniseysk and Krasnoyarsk. It is reported to have died out in the 18th century, and the Arin people were assimilated, today being associated with the Khakas. Only a small amount of data survives from early explorers' word lists. No speakers, schools or revitalization programme were found in the search results.

\section{Primary Educational Resources}
None found. No primer, textbook or school programme for Arin exists in the search results.

\section{Linguistic and Reference Materials}
\begin{itemize}
\item Wikipedia, Arin language: \url{https://en.wikipedia.org/wiki/Arin_language}
\item Wikipedia, Yeniseian languages: \url{https://en.wikipedia.org/wiki/Yeniseian_languages}
\item Wikipedia, Arin people: \url{https://en.wikipedia.org/wiki/Arin_people}
\item Wikipedia, Pumpokolic languages (a debated grouping that may include Arin): \url{https://en.wikipedia.org/wiki/Pumpokolic_languages}
\item Britannica, Arin language: \url{https://www.britannica.com/topic/Arin-language}
\item Sample pages of a scholarly volume from Reichert Verlag that surfaced in the search (content not verified): \url{https://medialibrary.reichert-verlag.de/media/pdf/9783752003215_sample.pdf}
\end{itemize}

\section{Religious and Community Resources}
None found.

\section{Digital and Audio Resources}
A popular blog post on the language's extinction: \url{https://www.freelang.net/blog/?p=802}. No audio resources exist, since the language has no recorded speakers.

\section{Curriculum Development Context}
Arin is attested only in fragmentary historical word lists, so a K-12 curriculum could not use the language as a medium of instruction. A heritage unit could place Arin within the history of the Yeniseian family and Siberian peoples, drawing on the reference works above. Related Yeniseian languages, such as Ket, would be the main source of material for any language-learning component.

\end{document}